% =============================================================================
% Which of the six Annual Review of Financial Economics surveys in Artigos/ is the most
% quant-focused? Companion to NotebookLM/audio/artigos-audio-part1..3.md.
% Compile: pdflatex (twice); check: pdffonts (Type 1 only, uni = yes).
% =============================================================================
\documentclass[11pt,a4paper]{article}
\usepackage[utf8]{inputenc}
\usepackage[T1]{fontenc}
\usepackage{lmodern}   % Type 1 vector fonts (replaces PK bitmaps)
\usepackage{cmap}      % ToUnicode CMap -> copyable text and accents
\usepackage[english]{babel}
\usepackage{amsmath,amssymb}
\usepackage[a4paper,margin=2.2cm]{geometry}
\usepackage{booktabs}
\usepackage{tabularx}
\usepackage{enumitem}
\usepackage{xcolor}
\usepackage{microtype}

\setlist{topsep=2pt,itemsep=2pt}
\newcolumntype{Y}{>{\raggedright\arraybackslash}X}
\newcommand{\verdict}[1]{\par\medskip\noindent\fbox{\parbox{\dimexpr\linewidth-2\fboxsep-2\fboxrule}{#1}}\par\medskip}

\begin{document}

\begin{center}
  {\Large\bfseries Which survey in \texttt{Artigos/} is the most quant-focused?}\\[4pt]
  {\large A ranking of six Annual Review of Financial Economics articles for a reader interested in
  quantitative trading (and green finance)}
\end{center}

\noindent\textbf{Articles.} Hong \& Shore, ``Corporate Social Responsibility'' (2023); Bolton, Chen \&
Wang, ``The Marginal Value of Cash'' (2024); Benmelech, ``The Benefits and Costs of Secured Debt''
(2024); Kermani \& Ma, ``Two Tales of Debt'' (2026); Colla, Ippolito \& Li, ``Debt Structure'' (2020);
Ewens \& Farre-Mensa, ``Private or Public Equity?'' (2022). References below are to each article's
sections, equations, tables and figures.

\verdict{\textbf{Bottom line.} \emph{The Marginal Value of Cash} (Bolton, Chen \& Wang) is the most
quant-focused in both senses of the word: it is the most mathematical (a solved dynamic model) and the
only one with a full asset-pricing section linking a firm characteristic to beta and expected returns.
\emph{Corporate Social Responsibility} (Hong \& Shore) is the runner-up and the strongest on empirical
return evidence. It is also the only article with green-finance content.}

\section*{1. What ``quant-focused'' means here}
Two different things, scored separately:
\begin{enumerate}[label=(\roman*)]
  \item \textbf{Quantitative modelling}: a formal model with equations, solved and calibrated.
  \item \textbf{Quant-trading relevance}: a priced, measurable relation (a return or a spread
        associated with a variable one could sort on in CRSP, Compustat, DealScan or ISS), a credible
        identification design, and warnings about the traps a backtest would hit (endogeneity, selection,
        omitted risk factors, short samples).
\end{enumerate}

\section*{2. The winner: \emph{The Marginal Value of Cash}}
\textbf{Modelling (§2--3).} The static model (Froot--Scharfstein--Stein; Kaplan--Zingales) maximises
$f(I)+w-I-C(e)$ subject to $e=\max(I-w,0)$. The first-order condition $f'(I)=1+C'(e)$ makes the marginal
value of cash $P_w=1+C'(e)\ge1$, which exceeds one exactly when the firm is constrained; the first best
($f'(I^{FB})=1$) is independent of $w$, which is Modigliani--Miller. Figure~4 calibrates it
($I^{FB}=0.30$, external finance used only for $w<0.23$). The dynamic model (§3, after Bolton, Chen \&
Wang 2011) is a continuous-time $q$-theory with frictions solved through an HJB equation: a payout
boundary fixed by a super-contact condition, an equity-issuance boundary, a credit-line limit and
hedging. Every policy depends on one state variable, $w=W/K$.

\medskip\noindent\textbf{Asset pricing (§4).} This is the part a quant can use:
\begin{itemize}[leftmargin=1.4em]
  \item \emph{Conditional beta} (Eq.~28): $\beta(w)=\dfrac{\rho\sigma}{\sigma_m}\,\dfrac{p'(w)}{p(w)}$.
        The marginal value of cash enters beta directly, so beta moves with financial constraints.
  \item \emph{The textbook unlevering formula fails} (§4.1, Eq.~30, Fig.~8a): treating cash as negative
        debt underestimates beta by 24\% (starting from the payout boundary) or overestimates it by
        32\% (starting from zero cash); at the boundary, beta can fall to 81\% of the unconstrained
        firm's.
  \item \emph{Idiosyncratic volatility is priced through beta} (§4.2, Fig.~8b), with a sign that depends
        on the cash level, against the usual ``diversifiable, hence irrelevant'' view.
  \item \emph{A sort variable with a published spread}: the top cash-to-assets decile beats the bottom by
        0.69\% a month, equal-weighted (Palazzo 2012).
  \item \emph{Conditional evidence}: cash-rich firms outperform cash-poor firms at the onset of the 2008
        crisis (Duchin, Ozbas \& Sensoy) and in early 2020 (Ramelli \& Wagner); cash is positively related
        to credit spreads (Acharya, Davydenko \& Strebulaev).
  \item \emph{Built-in warnings}: cash holdings are chosen, so the cash--return relation is endogenous and
        state-dependent; the spread is equal-weighted, so it may be concentrated in small caps.
\end{itemize}

\section*{3. The other five}
\begin{center}\small
\begin{tabularx}{\linewidth}{@{}lYl@{}}
  \toprule
  Article & Quant content & Verdict \\ \midrule
  \emph{Corporate Social Responsibility} &
  The most return-based evidence: Table~1 (cost-of-capital wedges; sin stocks about $+3\%$ a year, a
  $-1.5\%$ green-minus-brown wedge), Table~2 (portfolio performance), and the authors' own regression
  discontinuity on ISS shareholder votes, 2006--2021 (§3.2, Eqs.~2--3, Table~8), with Carhart
  four-factor abnormal returns. Theory is one equilibrium relation, $r_S-r_U=-m/q$ (Eq.~1); the greenium
  distribution (Fig.~4) has mean $-1.41\%$. &
  Strong runner-up \\[3pt]
  \emph{Secured Debt} &
  Credit pricing: raw secured spreads look 150--200 bp higher because riskier firms pledge collateral
  (Fig.~2); firm-by-time fixed effects (§3.2.2, Eq.~1) give a secured premium of about 2 bp for
  investment grade and 55 bp below it, rising with each rating downgrade (§5). &
  Best for a credit quant \\[3pt]
  \emph{Two Tales of Debt} &
  Within-firm, within-year comparisons of interest rates across debt types (§4.2), and debt--tangibility
  regressions (§4.3). Focus is macro and contract design. &
  Little to trade \\[3pt]
  \emph{Debt Structure} &
  Descriptive measures of how concentrated firms' borrowing is (HHI, Excl90; §4.1), with regressions
  explaining it (§4.4). &
  No return content \\[3pt]
  \emph{Private or Public Equity?} &
  Venture and IPO market trends and regulation (NSMIA, SOX, JOBS Act); the only return figure is the
  $-2\%$ average announcement return on seasoned equity offerings. &
  Least quant \\ \bottomrule
\end{tabularx}
\end{center}

\section*{4. Ranking}
\begin{center}\small
\begin{tabular}{@{}lccc@{}}
  \toprule
  Article & Modelling & Quant-trading relevance & Green finance \\ \midrule
  The Marginal Value of Cash & 1st & 1st & none \\
  Corporate Social Responsibility & 2nd & 2nd & \textbf{only direct coverage} \\
  The Benefits and Costs of Secured Debt & 3rd & 3rd (credit) & none \\
  Two Tales of Debt & 4th & 4th & none \\
  Debt Structure & 5th & 5th & none \\
  Private or Public Equity? & 6th & 6th & none \\ \bottomrule
\end{tabular}
\end{center}
\noindent Only \emph{The Marginal Value of Cash} has a solved model. The other modelling ranks are
judgement calls: Hong \& Shore have an equilibrium relation (Eq.~1) and an explicit regression
specification, Benmelech has testable predictions and a fixed-effects specification, and the last three
are descriptive and empirical. First place in both columns is the robust part of the ranking. ``None''
under green finance means the article has no environmental content; in the other five, ``green'' appears
only as an author's name (Greenwald) or a figure colour.

\section*{5. How to read them}
\begin{enumerate}[leftmargin=1.4em]
  \item \textbf{The Marginal Value of Cash}, §2 and §4 first (the beta formula and the cash evidence),
        then §3 for the model.
  \item \textbf{Corporate Social Responsibility}, §1.2 (Eq.~1), Tables~1--2, §3.2 and §4. It is the one
        to read if green finance matters, and a model of how to test a return anomaly.
  \item \textbf{The Benefits and Costs of Secured Debt}, §3.2 and §5, for credit spreads.
  \item The other three only for corporate-finance background (debt contracts, private markets).
\end{enumerate}

\noindent\textbf{Caveat.} The numbers above are as the articles report them; none was re-estimated here.
A trading strategy built on them would need its own out-of-sample test with point-in-time data.

\end{document}
